% 7_conclusion.tex

\section{Conclusion}
\label{sec:conclusion}

We have presented a visual analytics platform that integrates multi-criteria spatial suitability analysis, what-if impact simulation, and temporal scenario contextualisation for electric vehicle charging infrastructure planning. The system enables planners to interactively adjust criterion weights---through direct manipulation or structured AHP pairwise comparison---observe the resulting suitability rankings across approximately 118,000 H3 hexagonal cells, place candidate charger deployments, evaluate their estimated impact on energy delivery, carbon reduction, grid stress, and population coverage, and contextualise these impacts against Distribution Future Energy Scenarios demand projections.

The platform introduces a design pattern that extends the conventional GIS-MCDA workflow beyond site selection into impact evaluation and temporal stress testing. By overlaying scenario supply on DFES demand trajectories, the system answers a question that static suitability maps cannot: under different demand growth pathways, in which year will a proposed deployment become insufficient? This temporal dimension transforms planning from a one-time ranking exercise into an iterative process of scenario exploration and comparison.

The client-side architecture, built on DuckDB-WASM for analytical computation and H3 for spatial indexing, demonstrates that complex multi-criteria spatial analysis can be performed interactively in the browser without server infrastructure. This lowers the deployment barrier for planning teams and enables the system to function as a self-contained analytical tool that can be distributed alongside the datasets it analyses.

We demonstrated the platform through a case study of EVCP planning in Greater London, using nine heterogeneous open datasets. The use case walkthrough illustrated a complete planning cycle---weight configuration, spatial exploration via choropleth and glyph maps, multi-method sensitivity analysis, charger placement, impact estimation, DFES contextualisation, and multi-scenario comparison via radar charts. The results suggest that the interactive integration of these analytical stages enables a more transparent and exploratory approach to infrastructure planning than batch-oriented GIS-MCDA workflows.

Future work will address several directions. First, a formal user study with transport planners will evaluate the platform's usability and the decision quality it supports compared to existing workflows. Second, the impact model can be enhanced with time-varying assumptions (e.g., evolving grid emission factors and EV adoption rates) to improve the fidelity of temporal projections. Third, the modular architecture could be adapted to other spatial planning domains---such as renewable energy siting, healthcare facility placement, or urban greenspace allocation---by substituting the criteria layers, impact equations, and temporal reference data. Finally, collaborative features such as shared scenario libraries and annotation tools could extend the platform's utility for group decision-making processes~\cite{jankowski2001gis}.
